\section{Introduction}
\label{sec:intro}

In our dataset, a Sylius story about a disabled shipping method is not finished when it says the customer cannot place an order. The Gherkin file also has to capture the checkout state before the method is disabled, the channel/country condition, and the error shown at confirmation time. A Fineract loan story has a different shape: the important details often sit inside dates, repayment periods, and data tables. A Diaspora story can depend on visible browser behavior, such as profile navigation, tags, comments, or notifications. This is why Connextra-to-Gherkin generation is harder than rewriting one English sentence into another. The output must still read like a requirement, but it also has to look like a feature file that the project could actually keep. For BDD teams, the cost is not only writing the first draft. The cost also includes checking whether the generated steps match existing conventions and whether the scenario still expresses the original acceptance criteria.

Prior work already shows that LLMs are worth testing for BDD artifacts. Rathnayake et al. evaluated GPT-4, Claude 3, and Gemini, with GPT-4 reaching 91.16\% BERTScore F1 on their proprietary BDD dataset \cite{rathnayake2026bddgeneration}. That result supports the idea that large models can follow BDD-like text, but the private setting makes it hard to reuse as a public benchmark. Karpurapu et al. reported that few-shot prompting cut total Gherkin syntax errors by 89\% \cite{karpurapu2024bddautomation}. Their finding is important for our prompt design because it suggests that examples are not only decoration; they can change whether the model follows the Gherkin format. Ferreira et al. studied a GPT-4 Turbo acceptance-test pipeline in industry and reported 100\% syntactic correctness, but the Test Flow evaluation used 13 user stories from one company \cite{ferreira2025industrialcasestudy}. Fernandes et al. compared seven LLMs for Gherkin generation and reported METEOR 0.84 for Gemini; their sample had 10 test descriptions and did not include a BDD parser check \cite{fernandes2025gherkin}.

Those results do not answer our group scenario. The model is different: we evaluate \texttt{gpt-5.4-mini-2026-03-17}. The prompting is different: the example shown to the model comes from the same project family as the target story. The data are different as well. We use 261 Connextra-style stories reverse-engineered from open-source Gherkin: 100 from Sylius, 100 from Apache Fineract, and 61 from Diaspora. These projects do not share one writing style. Sylius scenarios are often about products, carts, channels, shipping methods, and tax rules. Fineract scenarios are more data-heavy and often use dates, loan balances, repayments, and tabular schedules. Diaspora scenarios use social-network and browser-navigation language. A benchmark that mixes these projects gives a stricter check than a single small backlog. It also tests whether same-domain few-shot prompting helps the model imitate project style rather than producing generic Given--When--Then text.

The study therefore tests two things at once. For semantic closeness, we reduce generated and reference files to Gherkin skeletons and compute cosine similarity with \texttt{all-MiniLM-L6-v2}. The skeleton step removes concrete names, numbers, and quoted values so that a product name or a date does not dominate the comparison. For parser readiness, we save the model output as a feature file and run \texttt{behave --dry-run}. We use a pinned GPT-5.4-mini version and \texttt{temperature=0} so that the run is repeatable. This paired design is important because a scenario can be close to the reference text and still fail the parser; the reverse is also possible when valid Gherkin leaves out a business condition. The two metrics therefore answer different practical questions: whether the generated scenario resembles the expert scenario, and whether a BDD tool can read it without syntax repair.

Section~\ref{sec:related} summarizes the related BDD and acceptance-test generation studies used to define the gap. Section~\ref{sec:method} describes the reverse-engineered dataset, prompt setup, API calls, and metric scripts. Section~\ref{sec:results} reports the Wilcoxon and binomial-test results. Section~\ref{sec:discussion} interprets the mixed outcome for BDD use. Section~\ref{sec:threats} gives the validity threats, and Section~\ref{sec:conclusion} answers the two research questions.
